\documentclass[11pt]{amsart}
\usepackage{amsmath,amssymb}
\usepackage[margin=1.1in]{geometry}
\usepackage{hyperref}
\newtheorem{theorem}{Theorem}
\newtheorem{lemma}[theorem]{Lemma}
\theoremstyle{remark}
\newtheorem{remark}[theorem]{Remark}

\title{Lov\'asz's path decomposition theorem from Fan's addible-edge lemma}
\author{Justin Williams}
\date{November 2026}

\begin{document}
\begin{abstract}
We give a short inductive derivation of Lov\'asz's theorem that a graph with at most one vertex of
even degree decomposes into $\lfloor n/2\rfloor$ paths, in the strong form in which every odd vertex is
an end of exactly one path. The only tool is Fan's addible-edge lemma. Both the derivation and Fan's
lemma are formally verified in Lean~4 with Mathlib.
\end{abstract}
\maketitle

Graphs are finite and simple. A \emph{path} is a nonempty simple path, and a \emph{decomposition} is a
partition of the edge set into paths. For a decomposition $D$ write $b_D(v)$ for the number of ends of
paths of $D$ at $v$, counted with multiplicity. Every path through $v$ uses two edges at $v$ if $v$ is
interior to it and one if $v$ is an end, so $b_D(v)\equiv d(v)\pmod 2$.

\begin{lemma}[Fan~{\cite[Lemma 3.4]{Fan}}]\label{lem:fan}
Let $uv\in E(G)$, let $G'=G-uv$, and let $D'$ be a decomposition of $G'$. If
\[
b_{D'}(v) > \bigl|\{w\in N_{G'}(u): b_{D'}(w)=0\}\bigr|,
\]
then $G$ has a decomposition $D$ with $|D|=|D'|$, $b_D(u)=b_{D'}(u)+1$, $b_D(v)=b_{D'}(v)-1$ and $b_D(w)=b_{D'}(w)$
for every other vertex $w$.
\end{lemma}

\begin{theorem}[Lov\'asz~\cite{L68}, strong form]\label{thm:lovasz}
If at most one vertex of $G$ has positive even degree, then $G$ has a decomposition $D$ with
$b_D(v)=d(v)\bmod 2$ for every vertex $v$. In particular $G$ decomposes into $|\{v:d(v)\text{ odd}\}|/2\le
\lfloor|V(G)|/2\rfloor$ paths.
\end{theorem}

\begin{proof}
We use strong induction on $\nu(G)=\sum_v d(v)+[\,G\text{ has no vertex of positive even degree}\,]$. Deleting an
edge $xy$ lowers the degrees of $x$ and $y$ by one and changes no other degree.

\emph{Case I: some vertex $z$ has positive even degree.} Choose a neighbour $u$ of $z$ and put $G'=G-uz$.
In $G'$ the vertex $z$ is odd. The vertex $u$ had odd degree in $G$, so in $G'$ it is even (possibly
isolated). Every other vertex keeps its parity. So $u$ is the only possible vertex of positive even
degree in $G'$, and $\nu(G')\le\nu(G)-1$. By induction $G'$ has a decomposition $D'$ with
$b_{D'}=d_{G'}\bmod2$; thus $b_{D'}(u)=0$ and $b_{D'}(z)=1$. Every neighbour of $u$ in $G'$ has positive
degree and differs from $u$, so it is odd and has an end. The right-hand side of Lemma~\ref{lem:fan} (with
the roles $u$, $v=z$) is therefore $0<1=b_{D'}(z)$. The lemma adds $uz$, moving an end from $z$ to $u$.
The result has $b(u)=1$ and $b(z)=0$, matching the parities in $G$.

\emph{Case II: every vertex of positive degree is odd, and some $x$ has $d(x)\ge3$.} Choose a neighbour $u$
of $x$, and let $H$ be $G-ux$ with a new leaf $\ell$ attached at $u$. In $H$ the vertex $u$ keeps its
degree, $\ell$ has degree $1$, and $x$ has positive even degree $d(x)-1\ge2$. So $x$ is the unique vertex of
positive even degree in $H$, and $\nu(H)=\nu(G)-1$ (same degree sum, indicator dropped). By induction $H$
has a decomposition $D_H$ with $b_{D_H}=d_H\bmod2$; in particular $b(x)=0$, $b(u)=1$, $b(\ell)=1$. The leaf is
an end of its path. Deleting $\ell$ gives a decomposition $D'$ of $G'=G-ux$ in which $b_{D'}(u)\in\{0,2\}$,
and all other end counts are unchanged. If the leaf path was the single edge $u\ell$, its only end at $u$
disappears; otherwise $u$ becomes an end of the shortened path.
\begin{itemize}
\item If $b_{D'}(u)=0$, add $ux$ as a new one-edge path.
\item If $b_{D'}(u)=2$, every neighbour of $x$ in $G'$ has an end: $u$ has two, and the others are odd. Lemma~\ref{lem:fan}
(with roles $x$, $v=u$) adds $ux$, moving an end from $u$ to $x$.
\end{itemize}
In both cases the result has $b(u)=b(x)=1$, and the number of paths is that of $D'$ or one more, as
required by the end count.

\emph{Case III: every vertex has degree at most $1$ and there is an edge $ux$.} Then $G-ux$ again has no
vertex of positive even degree and $\nu(G-ux)=\nu(G)-2$. Decompose it by induction and add $ux$ as a
one-edge path.

\emph{Case IV: $G$ has no edges.} The empty decomposition works.

The number of paths is $\tfrac12\sum_v b_D(v)$, which is half the number of odd vertices.
\end{proof}

\begin{remark}[Formal verification]
Theorem~\ref{thm:lovasz} is \texttt{lovasz\_ends} in \texttt{lean/GallaiLean/Lovasz.lean}. Lemma~\ref{lem:fan} is
\texttt{fan\_single} in \texttt{lean/GallaiLean/Fan.lean}, together with the grouped version
\texttt{fan\_grouped} (Fan's Lemma~3.3). Both are proved from the definitions in
\texttt{lean/GallaiLean/Basic.lean}: a decomposition is a list of simple walks of positive length whose
edge lists partition $E(G)$. The proofs use no \texttt{sorry}, and \texttt{\#print axioms} reports only
\texttt{propext}, \texttt{Classical.choice} and \texttt{Quot.sound}.
\end{remark}

\begin{remark}
We do not know whether this derivation of Lov\'asz's theorem from Fan's lemma appears in the
literature. Fan's lemma itself is proved by an exchange argument that does not use Lov\'asz's theorem, so
there is no circularity. Lov\'asz's original proof~\cite{L68} proceeds through the more general statement
on decompositions into paths and cycles.
\end{remark}

\begin{thebibliography}{9}
\bibitem{Fan} G.~Fan, Path decompositions and Gallai's conjecture, J.~Combin.\ Theory Ser.~B 93 (2005) 117--125.
\bibitem{L68} L.~Lov\'asz, On covering of graphs, in: Theory of Graphs (Proc.\ Colloq.\ Tihany, 1966),
Academic Press, 1968, 231--236.
\end{thebibliography}
\end{document}
